\chapter{Lernziele Kryptologie}

\begin{todolist}
\item Ich weiss wie die folgenden Kryptosysteme funktionieren (Verschlüsslung und Entschlüsslung bei gegebenem Schlüssel):
  \begin{todolist}
  \item Skytale 
  \item Caesar
  \item Vigenère
  \item One-Time-Pad
  \item RSA
  \end{todolist}
\item Ich kann die Verschlüsslung und Entschlüsslung mit Caesar und Vigenère in Python implementieren.
\item Ich kann einen Kryptotext, der durch die Vigenère-Verschlüsslung entstanden ist, mit Hilfe des Tools auf folgender Webseite entschlüsseln: \\
\url{https://cryptbreaker.marcwidmer.xyz}
\item Ich kann die Idee der Häufigkeitsanalyse erklären.
\item Ich kann eine gruppenweise Häufigkeitsanalyse durchführen, um Vigenère bei bekannter Schlüssellänge zu knacken.
\item Ich kenne das Kerckhoff’sche Prinzip der Sicherheit.
\item Ich kann die Friedman’sche Charakteristik für einen kurzen Text von Hand berechnen
\item Ich kann mit der Friedman'schen Charakteristik die Länge eines Vigenère-Schlüssels bestimmen.
\item Ich kann erklären, was mit der Friedman’sche Charakteristik berechnet wird.
%\item Ich verstehe das Beispiel des \textit{public-key-}Verfahrens mit dem Graphen (Schlüssel: perfect dominating set) und kann erklären, warum hier kein Schlüsselaustausch zwischen Bob und Alice notwendig ist.
%\item Ich verstehe das Prinzip hinter public Key Verfahren.
\item Ich kann eine Nachricht mit RSA asymmetrisch ver- und entschlüsseln, indem ich ein Beispiel mit kleinen Zahlen durchführe.
\end{todolist}
